\documentclass[10pt,a4paper]{amsart}
\usepackage[margin=19mm]{geometry}
\usepackage{amsmath,amssymb,mathtools}
\usepackage[scr=rsfs,cal=euler]{mathalfa}
\usepackage{xfrac}
\usepackage[unicode,hidelinks,bookmarks=false]{hyperref}
\newcommand{\ie}{i.e.\ }
\newcommand{\cf}{cf.\ }
\newcommand{\resp}{respectively\ }
\newcommand{\Subsection}{\subsection}
\providecommand{\nobreakdash}{\nobreak\hbox{-}}
\setlength{\emergencystretch}{2em}
\multlinegap0pt
\usepackage{tikz}
\usepackage{amssymb}
\usepackage{bm}
\usepackage{enumerate}
\usepackage{mathtools}
\usepackage{xcolor}
\usepackage{cancel}
\newtheorem{theorem}{Theorem}
\newtheorem{proposition}{Proposition}
\usepackage{cleveref}
\crefname{theorem}{Theorem}{Theorems}
\crefname{proposition}{Proposition}{Propositions}
\newcommand{\altTam}[2]{\operatorname{Tam}_{#1}(#2)}
\newcommand{\down}{^{\operatorname{down}}}
\DeclareMathOperator{\row}{Row}
\DeclareMathOperator{\udeg}{udeg}
\newcommand{\up}{^{\operatorname{up}}}
\title{Invariance of rowmotion \\ for variants of the Tamari lattice}
\author{Ben Adenbaum, Emily Barnard, Cesar Ceballos, Cl\'{e}ment Chenevi\`{e}re, Colin Defant, Sam Hopkins, Matthias M\"{u}ller, Martin Rubey, Jessica Striker}
\date{Source: arXiv:2609.12983v1 (CC BY 4.0)}
\begin{document}
\maketitle

\noindent\emph{Source and licence.} Invariance of rowmotion \\ for variants of the Tamari lattice. Version arXiv:2609.12983v1 (CC BY 4.0), distributed by arXiv under the Creative Commons Attribution 4.0 International licence, http://creativecommons.org/licenses/by/4.0/, as recorded in the arXiv licence metadata for this exact version and shown on the abstract page https://arxiv.org/abs/2609.12983v1. Full licence notice: \texttt{licenses/arxiv-2609.12983-CC-BY-4.0.txt}.\par\smallskip

\noindent\textit{Editorial scope.} Counted unit: the article's theorem thm:varphi\_udeg (source0.tex lines 3982--3986). The article's complete original proof of it is delivered with it (source0.tex lines 4056--4069, one \texttt{proof} environment of 14 lines, which the article places away from the statement). No mathematics is written, repaired or re-derived: the statement and the proof are the article's own lines, in the article's own notation, and no retained line is substituted or re-flowed.  Retained uncounted as the source-local context the proof needs: lines 3598--3600; lines 4037--4040.  Nothing was omitted from any retained span, so the package carries no omission record.  No retained line contains a citation, so the wrapper carries no bibliography and no bibliography entry.  No retained line contains a figure environment, an image-inclusion command or drawing code, so no image is delivered and the package declares no asset dependency.  Editorial formatting only, in addition: an AMS \texttt{amsart} preamble that restates the article's own declarations and macros used by the retained lines, namely the environments \texttt{theorem}, \texttt{proposition} on separate counters, together with the standard packages those lines need.  The retained environments print in this document as Theorem 1, Proposition 1 in order of appearance, whereas the article prints the same items with the numbering of its own full document.  The complete native source of the article as distributed in the arXiv e-print for this exact version is bundled unchanged as \texttt{sources/210166/source0.tex}.
\begin{theorem} \label{thm:row_varphi_commute}
We have $\widetilde \row \circ \varphi = \varphi \circ \row$.
\end{theorem}

\begin{proposition} \label{prop:udeg_split}
For every $(\delta,\nu)$-tree $T$, the map $\varphi$ preserves the lower up-degree, while the upper up-degree is transformed according to rowmotion:
\[\udeg\down(\varphi(T)) = \udeg\down(T); \qquad \udeg\up(\varphi(T)) = \udeg\up(\row(T)).\]
\end{proposition}

\begin{theorem} \label{thm:varphi_udeg}
The bijection $\varphi\colon \altTam{\nu}{\delta} \to \altTam{\nu}{\widetilde{\delta}}$ preserves orbit averages of the up-degree statistic. More precisely,
\[ \sum_{y \in \{\row^j(x)\colon j\ge 0\}} \udeg(y) = \sum_{y \in \{\widetilde{\row}^j(\varphi(x))\colon j\ge 0\}} \udeg(y) \]
for every $x \in \altTam{\nu}{\delta}$.
\end{theorem}

\begin{proof}[Proof of \cref{thm:varphi_udeg}]
Since $\varphi$ intertwines rowmotion by \cref{thm:row_varphi_commute}, we have
\[\sum_{y \in \{\widetilde{\row}^j(\varphi(x))\colon j\ge 0\}} \udeg(y) = \sum_{y \in \{\row^j(x)\colon j\ge 0\}} \udeg(\varphi(y)).\]
Moreover, 
\begin{align*}
\sum_{y \in \{\row^j(x)\colon j\ge 0\}} \udeg(\varphi(y)) &= \sum_{y \in \{\row^j(x)\colon j\ge 0\}} \udeg\down(\varphi(y)) + \sum_{y \in \{\row^j(x)\colon j\ge 0\}} \udeg\up(\varphi(y)) \\
&= \sum_{y \in \{\row^j(x)\colon j\ge 0\}} \udeg\down(y) + \sum_{y \in \{\row^j(x)\colon j\ge 0\}} \udeg\up(\row(y)) \\
&= \sum_{y \in \{\row^j(x)\colon j\ge 0\}} \udeg\down(y) + \sum_{y \in \{\row^j(x)\colon j\ge 0\}} \udeg\up(y) \\
&= \sum_{y \in \{\row^j(x)\colon j\ge 0\}} \udeg(y),
\end{align*}
where in the second equality we used \cref{prop:udeg_split}, and in the third equality we reindexed the orbit of rowmotion for the second term. Combining the two displayed identities yields
\[\sum_{y \in \{\widetilde{\row}^j(\varphi(x))\colon j\ge 0\}} \udeg(y) = \sum_{y \in \{\row^j(x)\colon j\ge 0\}} \udeg(y), \]
as required.
\end{proof}

\end{document}
